\section{Introduction}\label{intro:sec:introduction}

The discrete Fourier transform is a basic test case for arithmetic
complexity. Its familiar $O(n\log n)$ algorithms make the corresponding
lower-bound question particularly sharp: does every exact computation
require a constant multiple of $n\log n$ operations? The answer depends
on the computational model. We consider complex linear circuits with
arbitrary prechosen coefficients and prove that this lower bound is
false in its unrestricted nonuniform form.

The argument centers on a finite question about tensor products.
For integers $q,b\ge2$, applying a fixed $q\times q$ matrix on each axis
computes its $b$th tensor power with $bq^{b-1}$ matrix calls.
A \emph{monomial matrix} is a permutation matrix times an invertible
diagonal matrix. We prove that, for at least one invertible nonmonomial
matrix, an exact computation improves this
count. A single such finite saving is enough: tensor amplification
turns it into an exponent improvement, and a factorization that preserves
the number of coordinates transfers the improvement to Fourier matrices.
The existence proof develops a second set of tools, based on prices of
matrix calls and their behavior under feedback and changes of boundary
conditions.

\subsection{The model and the result}

For $n\ge2$, put $\zeta_n=\exp(2\pi i/n)$ and let
\begin{equation}\label{intro:eq:fourier}
 F_n=(\zeta_n^{jk})_{0\le j,k<n}.
\end{equation}
A linear circuit has inputs $x_0,\ldots,x_{n-1},0$ and a finite directed
acyclic graph of scalar gates. Each gate computes $u+v$, $u-v$, or
$\lambda u$ from previously available values, with a fixed
$\lambda\in\C$. Each such gate costs one operation, including every
scalar multiplication. Previously computed values can be reused
arbitrarily, and the outputs are designated available values. A fixed
combination $\alpha u+\beta v$ therefore costs at most three gates.

The graph and coefficients may depend on $n$, but not on the input.
Coefficients have no bound on magnitude, algebraic degree, or description
length; their generation is not charged. Depth, storage, cancellation,
and conditioning are unrestricted. All $n$ output coordinates must equal
$F_nx$ exactly for every $x\in\C^n$. Let $L(n)$ be the minimum number of
gates in such a circuit.

\Needspace{8\baselineskip}
\begin{theorem}\label{thm:main}
In this model,
\[
 \liminf_{n\to\infty}\frac{L(n)}{n\log_2 n}=0.
\]
Equivalently, for every $c>0$ and every integer $N_0\ge2$, there is an
integer $n\ge N_0$ and an exact circuit for $F_n$ with fewer than
$c n\log_2 n$ gates.
\end{theorem}

The $\Omega(n\log n)$ Fourier lower-bound question is discussed in
\citet[Abstract]{ailon2013} and \citet[Section~1]{almanRao2023}.
In the model above, \cref{thm:main} refutes the assertion that
$L(n)\ge c n\log_2 n$ for some fixed $c>0$ and all sufficiently large $n$.
The proof gives an unbounded sequence of
lengths, formed as products of distinct primes, with a vanishing
normalized gate count. It is an existence theorem in the stated
nonuniform model. The coefficients are chosen separately for each
resulting circuit; their preparation is outside this model. The theorem
concerns arithmetic operation counts along those lengths and makes no
assertion about coefficient size, numerical stability, or bit complexity.
Neither a uniform construction nor an all-length upper bound is needed
for this quantified conclusion.

\subsection{Historical context}

The fast Fourier transform of Cooley and Tukey organizes a composite-length
transform into smaller transforms and diagonal factors
\citep{cooleyTukey1965}. Coprime factors have a particularly useful
structure: Chinese remainder permutations identify the transform with a
tensor product without intervening diagonal factors. This is the
Good--Thomas factorization; Good's treatment gives the multidimensional
and coprime formulas, and Cooley, Lewis, and Welch describe Thomas's
construction and its relation to Good's work
\citep[Sections~11--12]{good1958}
\citep[pp.~1675--1677]{cooleyLewisWelch1967}.
Convolution supplies another classical route to general lengths.
The 1969 chirp-transform paper of Rabiner, Schafer, and Rader derives
the reduction to diagonal scalings and a convolution, crediting
Bluestein's 1968 work for the quadratic-exponent identity
\citep[Section~II]{rabinerSchaferRader1969}. Bluestein's journal treatment
appeared in 1970 \citep{bluestein1970}.
The transfer proved here uses the coprime tensor structure together with
a local factorization on exactly the original coordinates. Preserving
that width is what permits a finite tensor saving to survive synchronization.

Lower bounds have been proved under substantial restrictions on the
algorithm. Morgenstern's determinant argument establishes a bound of
order $n\log n$ for the unnormalized transform in a bounded-coefficient
model \citep{morgenstern1973}. Ailon proves a lower bound for the
normalized transform in a model with exactly $n$ registers and unitary
$2\times2$ gates \citep[Section~2 and Theorem~2.1]{ailon2013} and, in a
separate model, relates the operation count to the conditioning of
every intermediate transformation
\citep[Definition~2.1 and Theorem~3.1]{ailon2014}.
Those hypotheses are significant here:
the circuits in \cref{thm:main} are allowed arbitrary coefficients and
arbitrary intermediate conditioning.
An unrestricted scalar circuit can pass between normalized and
unnormalized transforms with $O(n)$ additional gates. Such a scaling
need not respect a prescribed coefficient bound, so the normalization
must be retained when comparing the restricted models.

Upper bounds have also improved within the $n\log n$ scale. In particular,
Alman and Rao obtain new leading constants for power-of-two transforms,
measured in real arithmetic operations \citep[Theorem~1.2]{almanRao2023}.
A fixed leading-constant improvement
does not settle the quantified question addressed here. The essential
step in this paper is to turn one strict finite tensor saving into a
smaller exponent in the tensor order, and then preserve that saving
when passing to an unbounded sequence of Fourier transforms.

Tensor amplification and duality also appear in recent work on
Kronecker-power circuits. Alman and Li use asymptotic-spectrum duality
to characterize the asymptotic size of depth-two linear circuits
\citep[Sections~1.3 and~4]{almanLi2025}. Their circuits are represented by rank-one
decompositions, with size measured by the nonzero entries in the two
matrix factors. Our finite-win question concerns sequential invertible
matrix calls on a fixed coordinate set. Its price construction must
therefore preserve the order of each program while combining different
matrix types.

The local Fourier construction uses established ideas from structured
linear algebra and reversible computation. Kuznetsov gives an
$LDL^{\mathsf T}$ factorization of geometric Vandermonde matrices,
including Fourier matrices, whose triangular factors are diagonal
scalings of triangular Toeplitz matrices
\citep[Theorem~1]{kuznetsov2018}. We give a Newton-polynomial derivation
of this identity and implement its factors on the original coordinates.
Low displacement rank supplies the convolution representation used in
that implementation \citep[Lemmas~1--2]{kailathKungMorf1979}.
Computing, reading the output, and reversing a computation is the
uncomputation pattern of Bennett \citep[pp.~526--529]{bennett1973}.
Restoration from arbitrary workspace contents is achieved by the
cleanup construction for transparent programs
\citep[Section~3]{buhrmanEtAl2014}. We prove the particular linear
replay identity, its pair-layer bound, and its fit inside the existing
coordinate set.

The companion manuscript \emph{An explicit power saving for the exact
discrete Fourier transform} gives a quantitative all-length algorithm
in a model that also charges scalar preparation and logarithmic-word
address operations, with a specified root of unity supplied
\citep[Theorem~1.1]{openaiExplicitFourier2026}.
The present paper develops the general finite-win existence argument:
comparison packing, the resulting price laws, and the reflection
contradiction are proved in full independently of that algorithm.

\subsection{The mechanism of the proof}

For an invertible $q\times q$ matrix $A$, the usual computation of
$A^{\otimes b}$ uses $bq^{b-1}$ calls to $A$, acting along tensor fibers.
Call a strict improvement a \emph{finite win}: an exact word on $q^b$
coordinates, using fewer calls, with permutations and invertible
diagonal maps allowed between them. The matrix $A$ must be nonmonomial.
This intermediate convention is confined to matrix words. Every scalar
multiplication is charged when such a word is implemented as a circuit.

The argument has two parts.

\begin{enumerate}[leftmargin=*,label=\textup{(\roman*)}]
\item \emph{A finite win implies \cref{thm:main}.}
For any nonmonomial $A$, one fixed positive-call pattern, with adjustable
monomial factors, realizes every invertible matrix of the same width.
A finite win then gives circuits for $A^{\otimes k}$ with
$O(q^k(k+1)^\alpha)$ gates for some fixed $\alpha<1$.
Each Fourier matrix $F_r$ is factored, on its original $r$ coordinates,
into $O(1+\log^4 r)$ layers, each monomial or a disjoint union of
invertible two-coordinate maps with any remaining coordinates fixed.
The temporary convolution computations
borrow arbitrary existing values and restore them by cancellation.
Positive generation places the pair maps in fixed-width tuples and
replaces them by a common call pattern in $A$. Synchronizing this pattern
across distinct prime factors makes each call slot a direct sum of
tensor powers of $A$. The sector widths sum to the Fourier length, so
the amplified bound controls the entire slot. The Chinese Remainder
Theorem then gives the desired sequence.

\item \emph{A finite win must exist.}
Suppose that none exists. Compare tensor-axis computations with words
using a finite collection of matrix types. If several comparisons
together save calls to every type, we can place their programs on
disjoint tensor sectors and schedule them as calls to one larger
block-diagonal matrix. This would give a finite win. Separation and
compactness turn the impossibility of such simultaneous savings into
one price function $p$, positive on nonmonomial matrices, with exact tensor
and direct-sum laws.

We next ask how prices change when coordinates are eliminated. Split
the coordinates of an invertible matrix into data and state, and apply
it repeatedly to fresh data while retaining the state. Repeat this
construction on many tensor axes. The cost of correcting the initial
state becomes negligible compared with the number of cell calls. This
proves that an invertible corner costs no more than its
containing matrix. Comparing multiplication modulo different scalar
polynomials also permits closing the state connection and specializing
rational matrix families. All these conclusions concern exact matrices;
they do not assume continuity of the price.

Finally, normalize a nontrivial elementary shear to have price one.
A legal exchange of one input basis member with one output basis member
preserves price, provided both exchanged lists remain bases. Applying
these exchanges to signed bit-coordinate matrices gives a family of
involutions, with integer parameter $d$, width $w=w(d)$, and price $dw$.
The constructed $+1$ eigenspace is the graph of an invertible matrix.
From the orthogonal reflection $R$ in that graph, we can recover the
involution's price at an additional cost of only $O(w)$.
The same subspace is the kernel of a real symmetric matrix with at most
$(d/2+1)w$ supported off-diagonal pairs. A product of diagonal and
two-coordinate perturbations has this sparse matrix as its derivative
at zero. At generic parameter values, each dense pair costs at most
$3/2$. Feedback adds only $O(w)$, and two algebraic specializations
carry the bound to $R$. Thus the two descriptions give
\[
 dw-O(w)\leq p(R)\leq \frac{3d}{4}w+O(w),
\]
with absolute implied constants. They are incompatible for sufficiently
large $d$.
\end{enumerate}

These parts meet at one precise interface: the finite word of
\cref{def:finite-win}. The transfer in \cref{prop:win-to-fourier} works
for any word satisfying that definition, while
\cref{thm:finite-win} establishes that at least one exists. In particular,
the proof does not require a finite win for every nonmonomial matrix.

\subsection{Conventions and order of the argument}

All matrices and scalar coefficients are over $\C$, except where real
symmetry and real orthogonality are explicitly used. Tensor products are
ordinary Kronecker products. For an invertible square
matrix $A$, its width is denoted by $|A|$. A matrix call acts on an ordered
tuple of distinct coordinates and fixes every other coordinate.
Zero-dimensional blocks may be omitted. A superscript $\mathsf T$
denotes ordinary transpose, not conjugate transpose.

The price function is introduced only under the assumption that there
is no finite win. It is an auxiliary real-valued function, not an
asserted lower bound for circuit size. Every limiting argument proves
an inequality for exact finite matrices. In particular, specialization
never replaces an exact circuit by a limit of approximate computations.

Constants in an $O$ bound may depend on fixed local matrices and their
dimensions. When a tensor order $k$ and a run length $t$ vary, each
polynomial bound in $k$ has a fixed exponent and is independent of $t$
unless stated otherwise. Scalar coefficients may depend on these
parameters, as the model permits. We specify the order of choices at
the packing, feedback, and specialization steps.

\Cref{gen:sec:generation,fft:transfer} prove the transfer from one finite
win, through positive generation, amplification, and exact-width Fourier
layers. \Cref{price:section} constructs the price function from the
opposite assumption. \Cref{corner:sec:corners,feedback:sec} establish
corner contraction, feedback, and algebraic specialization.
\Cref{pivot:sec:pivots} gives scalar-pivot invariance and the dense-pair
bound. Finally, \cref{reflect:section} uses those laws to force the
reflection contradiction and prove \cref{thm:finite-win}, which closes
the proof of \cref{thm:main}.

